% Section 6: what everything costs, measured and derived, and where the two
% algorithms sit against the literature. Sources are marked peer-reviewed or
% practitioner, because the maze-generation literature is thin and it matters.
\section{Time analysis and comparison}
\label{sec:analysis}

Write $n=WH$ for the number of cells.

\subsection{Size of the two structures}

\begin{center}
\begin{tabular}{@{}llr@{}}
  \toprule
  Output & What & How many \\
  \midrule
  Graph & nodes                   & $n$ \\
        & neighbour entries       & $2(n-1)$ \\
        & index $\langle P_i,P_j\rangle$ & $n$ \\
  Tree  & nodes                   & $2L-1$, for $L$ corridor leaves \\
  \bottomrule
\end{tabular}
\end{center}

Both are $\Theta(n)$, and the tree is the smaller of the two: $L$ measured
between $54$ and $8843$ on areas of $256$ to $65\,536$ cells, so well under $n$.
Keeping the second output costs a constant factor of memory and no asymptotics.

\subsection{Cost of generation}

Each call of \texttt{Divide} does $O(1)$ work of its own; each call of
\texttt{SpanCorridor} does work proportional to its corridor. Every door is
opened exactly once, and there are $n-1$ of them.

\begin{keypoint}
Generation runs in \keyterm{$\Theta(n)$ time} and fills \emph{both} outputs
while doing it. This is optimal: any generator must write $n-1$ doors.
\end{keypoint}

The recursion is not balanced --- the wall position is uniform, not central, the
same skew that makes quicksort's worst case quadratic. Measured depth against a
balanced tree on the same leaves:

\begin{center}
\begin{tabular}{@{}lrrr@{}}
  \toprule
  Area & Corridor leaves & Mean depth & Depth if balanced \\
  \midrule
  $16\times16$   & $54$   & $11.5$ & $5.8$ \\
  $64\times64$   & $666$  & $21.4$ & $9.4$ \\
  $256\times256$ & $8843$ & $32.0$ & $13.1$ \\
  \bottomrule
\end{tabular}
\end{center}

About $2.4$ times a balanced tree at every size --- the figure a stack-depth
bound has to be written against.

\subsection{Cost of solving the graph}

The maze is a tree, so $|E|=n-1$. With a binary heap Dijkstra costs
$O((|V|+|E|)\log|V|)=O(n\log n)$~\cite[Ch.~22]{cormen2022introduction}, and the
$\log$ buys nothing: all weights are $1$, so it settles cells in depth order and
is breadth-first search wearing a priority queue.

\subsection{Cost of solving the tree}

At each division either both endpoints stay on one side, and half the region is
discarded, or they are separated and both halves are entered. Write $p$ for the
probability of the second. Each call doing $O(1)$ work,
\[
  T(m) \;=\; (1-p)\,T(m/2) + p\cdot 2\,T(m/2) + O(1) \;=\; (1+p)\,T(m/2)+O(1),
\]
so by the master theorem

\begin{keypoint}
\[ T(n)\;=\;\Theta\!\left(n^{\log_2(1+p)}\right). \]
\end{keypoint}

The recurrence has one parameter, so it makes one prediction. Measuring $p$
directly, $200$ mazes per size:

\begin{center}
\begin{tabular}{@{}rrrrrr@{}}
  \toprule
  $n$ & One-way & Two-way & $p$ & $\log_2(1+p)$ & Exponent fitted to the visits \\
  \midrule
  $256$   & $10$  & $19$  & $0.65$ & $0.72$ & $0.70$ \\
  $1024$  & $27$  & $44$  & $0.62$ & $0.69$ & $0.69$ \\
  $4096$  & $69$  & $111$ & $0.62$ & $0.69$ & $0.68$ \\
  $16384$ & $170$ & $266$ & $0.61$ & $0.69$ & $0.68$ \\
  $65536$ & $435$ & $670$ & $0.61$ & $0.68$ & $0.67$ \\
  \bottomrule
\end{tabular}
\end{center}

Prediction and measurement agree at every size. The two ends of the scale are
worth reading off: at $p=1$ the recurrence is $2T(n/2)$ and the cost is
$\Theta(n)$, no better than Dijkstra; at $p=0$ it is $T(n/2)$ and the cost is
$\Theta(\log n)$.

\begin{keypoint}
$\Theta(\log n)$ is not available. It would need the endpoints never to be
separated, which happens only when they share a corridor. $\log n$ is the cost
of \keyterm{one} root-to-leaf descent, and the search pays one per separation.
\end{keypoint}

\subsection{The comparison}

$200$ mazes per size, start at $(0,0)$, exit at the far corner, counting cells
settled by the graph search and nodes visited by the tree search.

\begin{center}
\begin{tabular}{@{}lrrrrrrr@{}}
  \toprule
  & & & \multicolumn{4}{c}{graph output: cells examined} & tree \\
  \cmidrule(lr){4-7}\cmidrule(lr){8-8}
  Size & Cells & Steps & Dijkstra & tie: M & A*: M & A*: E & nodes \\
  \midrule
  $16\times16$   & $256$   & $64$   & $222$   & $221$   & $175$   & $191$   & $\mathbf{49}$ \\
  $32\times32$   & $1024$  & $180$  & $873$   & $870$   & $738$   & $776$   & $\mathbf{116}$ \\
  $64\times64$   & $4096$  & $538$  & $3533$  & $3530$  & $3199$  & $3280$  & $\mathbf{293}$ \\
  $128\times128$ & $16384$ & $1564$ & $13939$ & $13934$ & $12983$ & $13199$ & $\mathbf{702}$ \\
  $256\times256$ & $65536$ & $4493$ & $56193$ & $56186$ & $53753$ & $54276$ & $\mathbf{1776}$ \\
  \bottomrule
\end{tabular}
\end{center}

All $1000$ mazes were checked to be spanning trees. All of them were.

% Measured visits against area, on log-log axes, so that a power law is a
% straight line and its exponent is the slope.
\begin{figure}[htbp]
  \centering
  \begin{tikzpicture}
    \begin{loglogaxis}[
      width=0.86\textwidth, height=6.6cm,
      xlabel={cells $n$}, ylabel={nodes visited},
      log basis x=2, log basis y=10,
      legend pos=north west, legend cell align=left,
      legend style={font=\scriptsize, draw=black!30},
      grid=both, grid style={black!12},
      tick label style={font=\scriptsize}, label style={font=\small},
    ]
      \addplot[mark=*, wallred, thick] coordinates
        {(256,228) (1024,903) (4096,3543) (16384,13996) (65536,54293)};
      \addlegendentry{Dijkstra \quad slope $1.00$}
      \addplot[mark=square*, black!55, thick] coordinates
        {(256,68) (1024,173) (4096,536) (16384,1477) (65536,4348)};
      \addlegendentry{length of the route $d$}
      \addplot[mark=triangle*, doorgreen, thick] coordinates
        {(256,53) (1024,115) (4096,293) (16384,703) (65536,1588)};
      \addlegendentry{region tree \quad slope $0.68$}
      \addplot[dashed, black!45, no marks, domain=256:65536, samples=2]
        {sqrt(x)};
      \addlegendentry{$\sqrt{n}$, for reference}
    \end{loglogaxis}
  \end{tikzpicture}
  \caption{Both axes logarithmic, so a power law is a straight line and the
  exponent is its slope. Dijkstra runs parallel to $n$; the region-tree search
  runs at slope $0.68$, below even the length of the route it returns, because
  one visit to a corridor leaf covers a whole straight run.}
  \label{fig:growth}
\end{figure}


Three conclusions, in order of size.

\begin{description}[leftmargin=*,style=nextline,itemsep=6pt]
  \item[The tree beats the graph by a factor of thirty.] $1776$ against
    $56\,193$ at $256\times256$, and the gap widens with $n$ because the
    exponents differ: $\Theta(n^{0.68})$ against $\Theta(n)$. It also wins on
    memory --- $O(\log n)$ of recursion stack against $\Theta(n)$ of arrays
    Dijkstra must allocate before it starts.
  \item[Heuristic tie-breaking does nothing.] $5$ cells saved in $56\,193$, and
    Manhattan and Euclidean are \emph{identical} in every row. Both are
    predicted: cells nearer the start than the exit are necessarily
    expanded~\cite{dechter1985generalized}, so only the last layer is
    negotiable, and both estimates put the exit first in it.
  \item[A* helps a little, and Manhattan beats Euclidean.] $53\,753$ against
    $56\,193$, a $4\%$ saving, with Manhattan ahead at every size as its
    dominance predicts. Still $82\%$ of the maze.
\end{description}

\subsection{Why the heuristic is weak here, and the tree is not}

The two results above have one explanation between them.

A* prunes \emph{rival routes} --- on an open grid, the many equal-length ways to
reach the same cell. A perfect maze is a tree and has none, so all that is left
for A* to prune is dead-end branches whose depth plus estimate exceeds the
answer. And the estimate itself is poor: Sturtevant et
al.~\cite{sturtevant2009memory} measure that on $512\times512$ maze maps the
geometric estimate averages $151$ where the true distance is $512$--$768$, and
that a memory-based heuristic expands over eleven times fewer nodes.

\begin{keypoint}
A straight line says little about a route that must wind. The tree search wins
because it does not estimate anything: it \keyterm{knows} where the walls are,
having built them.
\end{keypoint}

For completeness: the grid-pathfinding methods with the largest speedups, Jump
Point Search and the symmetry-breaking family~\cite{harabor2011online}, work by
collapsing equal-cost alternative routes. A perfect maze has no path symmetry at
all, so they have nothing to prune here. Standard benchmark sets include maze
maps for that reason~\cite{sturtevant2012benchmarks}.

\subsection{Where recursive division sits}

The maze-generation literature is small and largely practitioner-written; the
distinction is marked because it bears on how much weight each claim carries.

\paragraph{The vocabulary.}
What this report builds is a \keyterm{perfect maze} --- any two cells joined by a
unique path --- which is the same object as a spanning tree of the grid graph.
The term is Pullen's~\cite{pullen2026think} and is used verbatim in
peer-reviewed work~\cite{bellot2021how}. Recursive division is a
\keyterm{wall adder}: it places walls where most generators carve passages.

\paragraph{The alternatives.}
Randomised depth-first search, randomised Kruskal and Prim, Eller's algorithm,
Hunt-and-Kill, binary tree and sidewinder all produce spanning trees of the same
grid in $\Theta(n)$ or close to it. Two do something this one does not:
Aldous--Broder~\cite{aldous1990random,broder1989generating} and
Wilson's algorithm~\cite{wilson1996generating} sample a \keyterm{uniform}
spanning tree, every perfect maze equally likely. Wilson's runs in expected mean
hitting time rather than cover time, a $\Theta(\log n)$ improvement on a grid.

\paragraph{Recursive division is not uniform.}
Its first division leaves exactly one door crossing a full straight line, and
recursively so; a spanning tree admitting no such decomposition can never be
produced. The support is a strict subset, so the distribution is not uniform
before one even asks how it is weighted inside it.
Pullen~\cite{pullen2026think} classifies it accordingly and records its visible
signature: long straight walls crossing the interior, rectangular rooms, a
self-similar texture. An $n$-cell grid has on the order of $3.21^{\,n}$ spanning
trees, and recursive division reaches only those with a guillotine
decomposition.

\begin{keypoint}
Recursive division was chosen for being \keyterm{simple, linear, recursive ---
and for the second output it gives free}. It was not chosen for the quality of
its randomness, and should not be claimed for it.
\end{keypoint}

\paragraph{What is missing from the literature.}
No peer-reviewed analysis of the distribution recursive division induces appears
to exist. Bellot et al.~\cite{bellot2021how} include it but score only
difficulty; Gabrov\v{s}ek~\cite{gabrovsek2019analysis} measures six generators
and omits it. Measuring its degree distribution against the uniform spanning
tree would close that gap, and is the natural continuation of this work.
